\section{\sc Top 10 Articles Demonstrating Scientific Impact} 


Dr.~Tschopp's accomplishments span research excellence, professional service, and mentorship. This section highlights research contributions that exemplify his scientific impact. His work focuses on the relationships between chemistry, processing, microstructure, and mechanics in materials---enabling innovations in how materials are manufactured, designed, and engineered for specific applications.

\vspace*{0.1in}
\begin{enumerate}
\item \textbf{Tschopp, MA, Ramsay, CW, Askeland, DR}, \textbf{\textit{Mechanisms of Formation of Pyrolysis Defects in Aluminum Lost Foam Castings}}, Trans.~AFS, 108 (2000) 609-614.
\begin{list2}
\item Award-winning work that identified key mechanisms behind defect formation in aluminum lost foam castings. Funded by DOE, this research was presented to 100+ industry partners and earned a Best Paper Award at the 2001 AFS Convention. At GM Powertrain, Dr.~Tschopp applied these findings to develop casting and validation processes that were adopted at manufacturing plants in Saginaw, MI and Defiance, OH. These changes led to multi-million-dollar annual savings across a new family of inline engine blocks and heads.
\end{list2}

\vspace{0.05in}

\item \textbf{Tschopp, MA, Spearot, DE, McDowell, DL}, \textbf{\textit{Atomistic simulations of homogeneous dislocation nucleation in single crystal copper}}, Modelling and Simulation in Materials Science and Engineering, 15 (2007) 693-709.
\begin{list2}
\item Pioneering study that challenged conventional theory by demonstrating that stress normal to the slip plane plays a major role in dislocation nucleation---orthogonal to Schmid stress-based models for dislocation motion. This marked the beginning of a broader body of work on dislocation nucleation mechanisms using atomistic simulations.
\end{list2}
\vspace{0.05in}

\item	\textbf{Tschopp, MA, McDowell, DL}, \textbf{\textit{Structures and energies of $\Sigma$3 asymmetric tilt grain boundaries in Cu and Al}}, Philosophical Magazine, 87 (2007) 3147-3173.
\begin{list2}
\item This highly cited paper launched a series of studies on symmetric and asymmetric grain boundaries in FCC metals. It provided a comprehensive analysis of how local atomic structure correlates with grain boundary character and grain boundary energy, often reproducing experimental HRTEM structures. At the time of publication, this was one of the most detailed investigations on the influence of grain boundary character and plane, influencing models of mechanical behavior in polycrystalline materials.
\end{list2}
\vspace{0.05in}

\item \textbf{Tschopp, MA, McDowell, DL}, \textbf{\textit{Dislocation nucleation in $\Sigma$3 asymmetric tilt grain boundaries}}, International Journal of Plasticity, 24 (2008) 191-217.
\begin{list2}
\item This foundational paper provided atomistic insights into how grain boundary structure, local stress state, and atomic-level mechanisms control dislocation nucleation in nanocrystalline FCC metals. This work helped establish a mechanistic basis for understanding deformation in ultrafine-grained and nanocrystalline materials.
\end{list2}
\vspace{0.05in}

\item \textbf{Tschopp, MA, Spearot, DE, McDowell, DL}, \textbf{\textit{Influence of Grain Boundary Structure on Dislocation Nucleation in FCC Metals}}, Dislocations in Solids: A Tribute to FRN~Nabarro,  vol. 14, pp. 43-139 (2008), ed. JP~Hirth.
\begin{list2}
\item This invited 97-page book chapter synthesized results from NSF-funded graduate work and presented a unified view of dislocation nucleation mechanisms in single crystals and grain boundaries. I delivered a keynote talk on this at the International Conference on Fatigue and this work formed the foundation of a dissertation that received the Sigma Xi Best PhD Thesis award.
\end{list2}
\vspace{0.05in}

\item \textbf{Tschopp, MA, Bartha, BB, Porter, WJ, Murray, T, Fairchild, S}, \textbf{\textit{Microstructure - dependent local strain behavior in polycrystals through in situ scanning electron microscope tensile experiments}}, Metallurgical Transactions A, 40 (2009) 2363-2368.
\begin{list2}
\item This experimental study introduced a methodology for coupling crystallographic orientation with local strain fields using nano-patterned speckles and digital image correlation in an SEM environment. It demonstrated how grain boundaries influence heterogeneous deformation, contributing widely cited tools and insights to the field of microscale experimental mechanics.
\end{list2}
\vspace{0.05in}

\item \textbf{Tschopp, MA, Wilks, GB, Spowart, JE}, \textbf{\textit{Multi-Scale Characterization of Orthotropic Microstructures}}, Modelling and Simulation in Materials Science and Engineering, 16 (2008) 065009.
\begin{list2}
\item Selected as the journal’s Editor’s Choice, this paper introduced novel synthetic microstructure generation techniques and multiscale statistical characterizations of composite microstructures. It influenced further work on image-based quantification and materials informatics methods, particularly in the context of Integrated Computational Materials Engineering (ICME).
\end{list2}
\vspace{0.05in}

\item \textbf{Hossain, D, Tschopp, MA$^\ast$, Ward, DK, Bouvard, JL, Wang, P, Horstemeyer, MF}, \textbf{\textit{Molecular dynamics simulations of deformation mechanisms of amorphous polyethylene}}, Polymer, 51 (2010) 6071-6083 ($^\ast$Corresponding Author).
\begin{list2}
\item This DOD-funded study employed atomistic simulations to investigate the molecular mechanisms underlying tension and compression in amorphous polyethylene. As corresponding author, Dr. Tschopp conceptualized and led the research, overseeing simulation development, designing post-processing algorithms, and guiding the project's scientific direction. He supervised and collaborated with a postdoctoral researcher (first author) and completed major portions of the work following the researcher’s departure. This highly-cited publication has become a widely referenced study in polymer mechanics at the atomic scale.
\end{list2}
\vspace{0.05in}

\item \textbf{Tschopp, MA, Solanki, KN, Gao, F, Sun, X, Khaleel, M, Horstemeyer, MF}, \textbf{\textit{Probing grain boundary sink strength at the nanoscale: Energetics and length scales of vacancy and interstitial absorption by grain boundaries in $\alpha$-Fe}}, Physical Review B, 85 (2012) 064108.
\begin{list2}
\item Published in a leading physics journal, this study leveraged high-throughput atomistic simulations to quantify point defect absorption energies across $\approx$200 grain boundaries in $\alpha$-Fe. It revealed the preferential binding of interstitials over vacancies at grain boundaries and complemented contemporaneous findings published in Science around the same time, contributing to defect transport models in radiation environments.
\end{list2}
\vspace{0.05in}

\item \textbf{Tschopp, MA, Murdoch, HA, Kecskes, LJ, Darling, KA}, \textbf{\textit{Bulk Nanocrystalline Metals: Review of the Current State of the Art and Future Opportunities for Copper and Copper Alloys}}, JOM, 66 (2014) 1000-1019.
\begin{list2}
\item This collaborative review examined the potential of nanocrystalline copper alloys for bulk applications, highlighting strategies for stabilizing ultrafine grain structures through chemistry and processing. The review synthesized emerging trends and key breakthroughs at ARL, laying groundwork for materials-by-design approaches targeting grain boundary engineering.
\end{list2}
\vspace{0.05in}

\end{enumerate}
